% SHARD-ID: CA-20-GALILEAN-PROJECTIVE-MASS
% SHARD-TITLE: Projective Galilean Symmetry and the Mass Central Term
% SHARD-SUMMARY: Explains why the projective-representation layer is unavoidable for Galilean quantum symmetry.
% SHARD-SUMMARY: Derives the mass central coefficient from the canonical commutator [mX_a,P_b]=i m delta_ab on a common core.
% SHARD-SUMMARY: Records the source gap for a registered Bargmann mass-extension reference while keeping the local coefficient checked.
% SHARD-KEYWORDS: Galilean projective representation, mass central extension, Bargmann source gap, canonical commutator, Stone theorem

\section{Projective Galilean Symmetry and the Mass Central Term}
\label{sec:galilean-projective-mass}

\subsection{Status, Sources, and Dependencies}

\textbf{Status: sourced projective framework plus checked local derivation.}
Schottenloher supplies the representation-theoretic framework: projective
symmetries need not lift to the original group, but lift through a central
extension.  The coefficient \(m\delta_{ab}\) below is a local canonical
commutator derivation checked in Julia.  A local source specifically identifying
the Galilei mass central extension by its standard name has not yet been
registered.

Dependencies: CA-05, CA-18, CA-19, and CONVENTIONS.md (f), (i).

% Source: references/cft/Schottenloher2008/Schottenloher2008.md:1524 -- unitary representation.
% Source: references/cft/Schottenloher2008/Schottenloher2008.md:1528 -- projective representation.
% Source: references/cft/Schottenloher2008/Schottenloher2008.md:1548 -- lifting projective representations question.
% Source: references/cft/Schottenloher2008/Schottenloher2008.md:1554 -- no lift to G in general; central extension of the universal cover.
% Source: references/cft/Schottenloher2008/Schottenloher2008.md:1570 -- central extension by U(1).
% Source: references/cft/Schottenloher2008/Schottenloher2008.md:1618 -- topological central-extension lift.
% Source: references/cft/Schottenloher2008/Schottenloher2008.md:1622 -- quantization as a central extension by U(1).
% Source: references/cft/Schottenloher2008/Schottenloher2008.md:2289 -- Bargmann theorem hypothesis for lifting projective representations.
% Source: references/cft/Schottenloher2008/Schottenloher2008.md:2293 -- projective representation lift under Bargmann theorem hypothesis.
% Checked: test/runtests.jl, "Galilei mass central coefficient".

\subsection{Why This Belongs in the Grammar}

CA-05 made a conservative grammar choice: a projective symmetry is represented
by an honest unitary representation of a named central extension
\[
  1\longrightarrow U(1)\longrightarrow E\longrightarrow G\longrightarrow 1.
\]
For Galilean symmetry this is not a decorative detail.  The central parameter is
physically interpreted as mass in the canonical nonrelativistic model, and it
controls how boosts and momenta compose at the quantum level.

\subsection{Canonical One-Particle Calculation}

Work formally on a common invariant core where the position and momentum
operators obey
\[
  [X_a,P_b]=i\delta_{ab}I.
\]
For a particle of mass \(m\), the time-zero boost generator is represented by
\[
  G_a(0)=mX_a.
\]
Then
\[
\begin{aligned}
  [G_a(0),P_b]
    &= [mX_a,P_b]\\
    = m[X_a,P_b]\\
    = i\,m\,\delta_{ab}I.
\end{aligned}
\]
The Julia check records exactly the real coefficient \(m\delta_{ab}\).

\subsection{Compatibility with Time Evolution}

For the free Hamiltonian
\[
  H=\frac{1}{2m}\sum_b P_b^2,
\]
the same canonical commutator gives
\[
\begin{aligned}
  [H,mX_a]
  &= \frac{1}{2}\sum_b [P_b^2,X_a]\\
  = \frac{1}{2}\sum_b\bigl(P_b[P_b,X_a]+[P_b,X_a]P_b\bigr)\\
  &= \frac{1}{2}\bigl(-iP_a-iP_a\bigr)
  = -iP_a.
\end{aligned}
\]
Thus
\[
  i[H,G_a(0)]=P_a.
\]
This matches the CA-19 vector-field relation \([H,G_a]=P_a\) after applying the
self-adjoint-generator convention used for observable evolution.

\subsection{Central Extension Versus Vector Fields}

The unextended vector-field algebra had \([G_a,P_b]=0\).  The canonical quantum
calculation has instead
\[
  [G_a,P_b]=i\,m\,\delta_{ab}I.
\]
The additional term is central: it is a scalar multiple of the identity on a
fixed-mass sector.  This is the algebraic reason the compiler must attach a
central-extension witness to Galilean symmetry data rather than store only the
spacetime vector fields.

\subsection{Source Gap}

The local source collection includes Bargmann's theorem about lifting projective
representations under a Lie-algebra cohomology hypothesis, and it includes the
general central-extension theorem for projective representations.  It does not
yet include a registered source for the standard mass central extension of the
Galilei group.  Until that source is acquired, this shard uses the cautious name
\emph{mass central term} and cites the checked canonical calculation as its
local evidence.

\subsection{Compiler Payload}

A Galilean projective symmetry record must therefore include:
\[
\begin{array}{ll}
\text{base group/algebra} & \text{Galilei target from CA-19},\\
\text{central generator} & M\text{ or }I\text{ on a fixed-mass sector},\\
\text{central charge} & m\text{ or a mass operator},\\
\text{raw commutator witness} & [G_a,P_b]=i\delta_{ab}M,\\
\text{sign convention} & \text{conversion to }i[A,B]\text{ stated explicitly}.
\end{array}
\]
An OR of sectors with different masses can still be a Hilbert direct sum, but
the central element must then be represented as a block-diagonal operator rather
than a single scalar charge.
